\documentclass[11pt]{article}
\usepackage[margin=1in]{geometry}
\usepackage{amsmath,amssymb}
\usepackage{booktabs}
\usepackage{enumitem}
\usepackage{microtype}
\usepackage[hidelinks]{hyperref}

\setlist{nosep,leftmargin=*}
\newcommand{\bq}{\mathbf{q}}
\newcommand{\bp}{\mathbf{p}}

\title{\textbf{Conditional Generation under Weak Conditioning:\\[2pt]
EEG-Conditioned Hand Motion Generation as a Testbed}\\[6pt]
\large Independent Study Proposal --- Deep Learning}
\author{Zhizhe Zhang (zz5070) \quad Dixuan Yang (dy2397)\\
NYU Shanghai \\[4pt]
\small Advisor: Prof.\ Shengjie Wang}
\date{September 2026}

\begin{document}
\maketitle
\vspace{-2em}

\begin{abstract}
\noindent Conditional generative models---classifier-free guidance, adaptive normalization, cross-attention conditioning---were tuned on \emph{strong} conditions such as text prompts and class labels. How they behave when the condition is weak (low mutual information with the target, low SNR, unstable across domains) is largely untested. We study that question with EEG-conditioned hand-trajectory generation on WAY-EEG-GAL as the testbed. The contributions are (i) a structured, low-dimensional output parameterization separating path geometry from execution timing; (ii) \emph{condition-fidelity} diagnostics that measure whether a conditional generator actually uses its condition, which existing EEG generation work does not test; and (iii) a controlled experiment in which condition strength is an independent variable, so that EEG is one point on a curve rather than the sole evidence. A two-week validity gate precedes all generative modeling, because reported EEG trajectory-decoding correlations may be largely explained by task stereotypy and metric bias.
\end{abstract}

\section{Motivation}

EEG motor decoding splits into two regimes. Movement-type classification has been replicated for decades: it needs only a few bits, permits averaging over the trial, and has a hard chance level. Continuous trajectory regression instead reports Pearson correlations (PCC) of 0.4--0.7 whose provenance is rarely verified, for three reasons. \emph{Chance is not zero}: Kobler et al.~\cite{kobler2018} obtained $r\approx0.10$--$0.13$ by shuffling, and their control condition---subjects watching a replayed cursor with the hand at rest---still yielded $r\approx0.33$--$0.36$ for ``hand velocity'' versus $r\approx0.40$ during real execution. \emph{The metric is biased}: PCC is scale-invariant and systematically high for slow, sinusoid-like signals, which both low-pass EEG and hand velocity are; against a randomly re-paired baseline, Antelis et al.~\cite{antelis2013} found reported correlations were not above chance. \emph{The task is stereotyped}: the 3{,}936 WAY-EEG-GAL trials share nearly identical kinematics, so a model that ignores its input and emits the phase-conditioned mean scores well.

The conclusion is not that regression fails, but that its true performance is an \emph{unmeasured quantity}. The first and third reasons are one phenomenon: the condition is too weak, so the model fits the marginal distribution of the output. That is a known failure mode of conditional generation which the BCI literature does not name as such; conversely, the conditional-generation toolkit was tuned on strong conditions, and how it degrades as conditions weaken has not been studied systematically. EEG is an extreme but honest instance of the general problem.

\section{Weak Conditioning}

``Weak'' conflates three properties with different remedies (Table~\ref{tab:weak}); separating them turns an excuse into a diagnosis.

\begin{table}[h]
\centering\small
\begin{tabular}{@{}p{2.5cm}p{3.1cm}p{3.9cm}p{4.2cm}@{}}
\toprule
Dimension & Meaning & Algorithmically addressable & Consequence \\
\midrule
Low information & Small $I(c;x)$ & No (hard ceiling) & Only remedy is coarser output \\
Low SNR & Signal buried in artifact & Yes (representation, denoising) & Encoder design \\
Distribution shift & Cross-subject variability & Yes (adaptation, invariance) & Cross-subject generalization \\
\bottomrule
\end{tabular}
\caption{Three independent dimensions of weak conditioning.}
\label{tab:weak}
\end{table}

\vspace{-0.5em}
Three failure modes follow and are the objects of study: \emph{condition collapse} (the model emits the marginal, i.e.\ the template); \emph{guidance amplifying noise} (CFG scales the conditional--unconditional difference, which for a noisy condition is largely noise, so the text-to-image default $w\approx7.5$ cannot transfer); and a \emph{deformed diversity--fidelity trade-off} (sample diversity may reflect honest uncertainty or a failure to learn the condition, indistinguishable under standard metrics).

\section{Related Work}

PreMovNet~\cite{premovnet}, ESI-GAL~\cite{esigal} and M3T-Attention~\cite{m3t} regress hand kinematics from low-frequency EEG on WAY-EEG-GAL and all report PCC; the public M3T script splits by session within one subject, so its scores do not reflect unseen-subject performance. Among generative neighbors, JET~\cite{jet} generates EEG via flow matching, NeuroSonic~\cite{neurosonic} reconstructs speech from EEG via conditional flow matching, and FDA~\cite{fda} learns neural-to-behavior maps on invasive recordings. Several of their choices transfer: JET's structure-preserving constraints against mean-seeking, its concave optimum in constraint weight, and its finding that a degenerate zero-noise base distribution is catastrophic; NeuroSonic's $x$-prediction with a velocity-space loss~\cite{jit} and its preference for deterministic ODE sampling under noisy conditioning.

What neither does is test whether the model uses its condition. NeuroSonic has no shuffle control, no template baseline, no condition dropout and no guidance analysis, and its metrics (FAD, LSD, spectral convergence, DNSMOS) measure distributional similarity or per-sample quality---a model ignoring EEG and generating plausible speech would score well. Closing that gap is our central methodological contribution, and the diagnostics below apply directly to that work. No prior work generates hand motion from EEG, and none treats condition strength as an experimental variable; hand motion additionally offers a separable path/timing structure with no analogue in speech.

\section{Method}

\paragraph{Task.} Input: roughly 1\,s of scalp EEG preceding movement onset, and nothing else. Output: the full 3D wrist displacement sequence of the subsequent reach-to-grasp (1--2\,s), relative to the onset position, generated in one pass with no EEG inside the prediction horizon. This avoids a degeneracy of lagged sliding-window decoding: over 300--500\,ms the displacement is nearly a line segment, and path and timing parameters become collinear.

\paragraph{Structured output.} Following the space--time decoupling of arc-length dynamic movement primitives~\cite{dmp}, a trajectory is written as $\Delta\bp(u)=\bq(s(u))$, where $\bq$ is the spatial path parameterized by cumulative arc length (8 spline control points, 24 dims) and $s(\cdot)$ a monotone progress map from time to arc position (10 dims). Reconstruction is differentiable and reduces the output from ${\sim}300$ to ${\sim}34$ dimensions. Two effects are intended: lowering the information the model must extract from the condition, and removing the mechanism by which per-timestep regression penalizes temporal misalignment and thereby rewards averaged trajectories. The decomposition is not novel; the claim concerns its role as a generative output parameterization under weak conditioning, and its collinearity at the chosen horizon is tested explicitly (Gate~2).

\paragraph{Generator.} A conditional flow-matching model~\cite{fm} generates path and timing parameters jointly, under the $x$-prediction / velocity-loss formulation of JiT and NeuroSonic. The backbone is a small Transformer with AdaLN injection of the flow time; the base distribution is standard Gaussian (no template initialization, per JET); condition dropout at $p=0.1$ enables CFG analysis and an unconditional baseline. Three injection variants are compared: AdaLN, cross-attention, and joint-token concatenation. Following JET, the loss adds L1 displacement and velocity reconstruction, a weak bell-shaped velocity-profile prior, a soft jerk bound and endpoint error; each term is ablated and the smoothness weight swept, since stronger smoothness yields prettier trajectories precisely by manufacturing conditional means.

\paragraph{Condition-fidelity diagnostics.} Three model-agnostic measures of whether a generator uses its condition: (i) a \emph{cross-conditioning matrix} $D_{ij}=\mathrm{err}(G(c_i),x_j)$, where condition usage implies diagonal dominance, reported with a permutation $p$-value; (ii) the \emph{within- vs.\ between-condition variance ratio}, where equality indicates collapse; and (iii) the \emph{silhouette score} of samples clustered by condition, as in JET. These apply to any conditional generator, including the discriminative baselines. A component-wise ablation (EEG fed to the path branch only, the timing branch only, or neither) doubles as a collapse diagnostic and as a finding about which motor component EEG conditions; a DTW-based decomposition separates spatial from timing error for any model's output.

\section{Experimental Design}

\paragraph{Gate 1 (weeks 1--2): does EEG carry usable information?} Four models under identical splits: a phase-conditioned template that ignores its input; a shuffled-pairing model retrained 50--100 times for an empirical chance distribution; an EMG-conditioned positive control (if EMG cannot beat the template, the pipeline is broken); and a lightweight EEG regressor. The shuffled model should match the template---exceeding it indicates leakage---and dropping frontal channels should not collapse performance to chance, which would indicate ocular contamination. \emph{Pass criterion, fixed in advance:} the EEG model's mean displacement error on held-out subjects falls below the 5th percentile of the shuffle distribution.

\paragraph{Gate 2 (weeks 3--4).} Reconstruction error of the path/timing representation, an explicit collinearity test, and reproduction of M3T.

\paragraph{Core $2\times2$ (weeks 5--7).} With encoder and training budget fixed, deterministic regression and conditional flow matching are each run on the raw coordinate sequence and on the path/timing representation, separating the contribution of the representation from that of generative training. Weeks 8--10 add the fidelity diagnostics on all models, the component-wise and injection ablations, a CFG sweep $w\in\{0,0.5,1,1.5,2,3,5\}$, ODE vs.\ SDE sampling, and the constraint ablation.

\paragraph{Condition strength as an independent variable (weeks 8--11, parallel).} Starting from a strongly conditioned motion-generation setting (clean intent labels or demonstration data), the condition is degraded along the three dimensions of Table~\ref{tab:weak}: additive noise, an information bottleneck, and per-``subject'' distribution shift. Repeating the core experiments at each level plots the gain from structured output, the optimal CFG scale, the collapse threshold and the ODE--SDE gap as functions of condition strength, with EEG as the extreme point. This experiment is not data-limited and is the primary answer to ``why EEG?''.

\paragraph{Evaluation (weeks 11--12).} Leave-one-subject-out, reporting per-subject distributions rather than a grand average, stratified by trial variability to test JET's prediction that structural gains concentrate on high-variability trials. Primary metrics are mean 3D displacement error (mm), endpoint error, velocity error and velocity-peak timing; per-axis PCC and $R^2$ are reported only for comparison with the literature; generative quality uses a motion-domain distributional distance over kinematic features (not image FID) and the single-sample error distribution, with no best-of-$K$ selection. Splits are made by subject, then session, then trial, before any windowing; normalization statistics come from training data only; filtering is causal except when reproducing prior protocols, where it is labeled as such.

\section{Data, Feasibility, and Risks}

WAY-EEG-GAL~\cite{wayeeggal} provides 12 participants and 3{,}936 grasp-and-lift trials, 32-channel EEG at 500\,Hz with synchronized hand position, EMG and force; holding out 3--4 subjects leaves roughly 2{,}600--2{,}900 training pairs. This is small for generative modeling (JET: 130M parameters on 10{,}000+ sessions; NeuroSonic: 49{,}200 pairs from 48 subjects) and is treated as a hard constraint: model size stays in the low millions of parameters, results are per-subject LOSO distributions, and the collapse diagnostics are mandatory, because small data plus weak conditions maximize the incentive to fit the marginal. Windowing increases sample count but not information and must follow the split. Public code exists for M3T, JET and NeuroSonic.

\begin{itemize}
\item \emph{Gate 1 fails.} Reduce output granularity to direction plus speed profile (an order of magnitude less information), or complete the method on the controlled condition-strength setting with EEG as a documented failure case---a bounded algorithmic result, though no longer an EEG paper.
\item \emph{Gate 2 shows degeneracy.} Replace the decomposition with a direction sequence plus scalar speed profile, which does not degenerate on short segments.
\item \emph{Core grid does not beat baselines.} The fidelity diagnostics, component-wise ablation and empirical rules for CFG and injection under weak conditioning stand on their own.
\end{itemize}

\section{Questions for the Advisor}

\begin{enumerate}
\item Is the level of abstraction right---a claim about conditional generation under weak conditioning, with EEG as testbed---or should the work be anchored more in motion generation and control, giving the condition-strength experiment greater weight?
\item Whole-segment generation (1--2\,s, avoiding degeneracy) versus short-horizon prediction (300--500\,ms, closer to existing baselines)?
\item Scope: one dataset, one reproduced baseline, core experiments. If Gate 1 fails, reduce granularity or move to a controlled demonstration-data setting?
\end{enumerate}

\begin{thebibliography}{9}\small
\bibitem{antelis2013} J.~M. Antelis et al. On the usage of linear regression models to reconstruct limb kinematics from low frequency EEG signals. \emph{PLoS ONE}, 8(4):e61976, 2013.
\bibitem{kobler2018} R.~J. Kobler, A.~I. Sburlea, G.~R. M\"uller-Putz. Tuning characteristics of low-frequency EEG to positions and velocities in visuomotor and oculomotor tracking tasks. \emph{Sci.\ Rep.}, 8:17713, 2018.
\bibitem{wayeeggal} M.~D. Luciw, E.~Jarocka, B.~B. Edin. Multi-channel EEG recordings during 3,936 grasp and lift trials. \emph{Sci.\ Data}, 1:140047, 2014.
\bibitem{premovnet} A.~Jain, L.~Kumar. PreMovNet: Pre-movement EEG-based hand kinematics estimation for grasp-and-lift task. \emph{IEEE Sensors Letters}, 2022.
\bibitem{esigal} A.~Jain, L.~Kumar. ESI-GAL: EEG source imaging-based kinematics parameter estimation for grasp and lift task. arXiv:2406.11500, 2024.
\bibitem{m3t} M3T-Attention: multi-level multi-scale temporal attention for EEG-based kinematics decoding. 2026. \url{https://github.com/jjspp/M3T_Attention}.
\bibitem{jet} Y.~Wang, Y.~Ma, W.~Li, C.~You. Let EEG models learn EEG. \emph{ICML}, 2026.
\bibitem{neurosonic} W.~Gao et al. NeuroSonic: Conditional flow matching for EEG-to-speech reconstruction. \emph{MICCAI}, 2026.
\bibitem{fda} P.~Wang et al. Flow matching for few-trial neural adaptation with stable latent dynamics. \emph{ICML}, 2025.
\bibitem{jit} T.~Li, K.~He. Back to basics: Let denoising generative models denoise. arXiv:2511.13720, 2025.
\bibitem{fm} Y.~Lipman et al. Flow matching for generative modeling. \emph{ICLR}, 2023.
\bibitem{dmp} Arc-length dynamic movement primitives. \emph{Robotics and Autonomous Systems}, 2017.
\end{thebibliography}

\end{document}
